\chapter{Dystonia}

\section*{Definition}
Dystonia is a movement disorder characterized by sustained or intermittent muscle contractions causing abnormal, often repetitive, movements, postures, or both. It may be focal, segmental, multifocal, generalized, or hemidystonic, and is classified by body distribution, age of onset, and etiology.

\section*{Pathophysiology}
Dystonia results from dysfunction within the basal ganglia-thalamocortical motor circuits, particularly involving impaired inhibition at multiple levels of the motor system. Reduced GABAergic inhibition, aberrant sensorimotor integration, and abnormal synaptic plasticity lead to excessive co-contraction of agonist and antagonist muscles. Genetic forms involve mutations in genes affecting neurotransmitter synthesis, ion channel function, and cellular trafficking pathways.

\section*{Etiology}
\begin{itemize}
  \item Primary (idiopathic): DYT1 (TOR1A) mutation causing early-onset generalized dystonia, DYT6 (THAP1) causing craniocervical dystonia, other monogenic forms (DYT4, DYT11, DYT12)
  \item Focal dystonia syndromes: Cervical dystonia (spasmodic torticollis), blepharospasm, oromandibular dystonia, laryngeal dystonia (spasmodic dysphonia), writer's cramp, musician's dystonia
  \item Secondary (acquired) causes: Stroke, traumatic brain injury, cerebral palsy (hypoxic-ischemic insult), brain tumors, encephalitis, multiple sclerosis, Wilson disease, neurodegeneration with brain iron accumulation
  \item Drug-induced: Antipsychotics (tardive dystonia), metoclopramide, levodopa (dopamine dysregulation), anticonvulsants
  \item Psychogenic (functional): Onset often abrupt, incongruent with known organic patterns, and accompanied by other functional neurologic symptoms
\end{itemize}

\section*{Indications for Evaluation}
Seek care if:
\begin{itemize}
  \item Persistent abnormal postures or movements interfering with function or quality of life
  \item Progressive symptoms, spreading to new body regions
  \item Onset before age 20 (higher likelihood of genetic etiology)
  \item Associated neurologic signs (weakness, spasticity, cognitive decline, ophthalmologic findings)
  \item History of antipsychotic or antiemetic use (suspect tardive dystonia)
\end{itemize}

\section*{Related Symptoms}
\begin{itemize}
  \item Tremor (dystonic tremor)
  \item Myoclonus
  \item Spasticity
  \item Chorea
  \item Parkinsonism
  \item Muscle pain
  \item Fatigue
  \item Anxiety
\end{itemize}
\textbf{Note:} Dystonia is often task-specific (e.g., writer's cramp occurring only during writing) and may be alleviated by sensory tricks (geste antagoniste) such as touching the affected body part.

\section*{Self-Care / Home Management}
\begin{itemize}
  \item Identify and utilize sensory tricks to temporarily reduce dystonic postures (e.g., gently touching the chin for cervical dystonia).
  \item Reduce stress, fatigue, and anxiety which commonly exacerbate dystonic movements.
  \item Use supportive devices such as padded chairs, ergonomic writing tools, or neck supports for symptomatic relief.
  \item Physical and occupational therapy focusing on stretching, relaxation, and adaptive strategies.
\end{itemize}

\section*{Diagnostic Approach}
\begin{itemize}
  \item Clinical diagnosis based on characteristic phenomenology: sustained muscle contractions, patterned movements, twisting postures, and geste antagoniste.
  \item Family history and age of onset guide genetic testing (DYT1, DYT6, DYT11, and others).
  \item Brain MRI excludes secondary structural causes (stroke, tumor, Wilson disease, neurodegeneration).
  \item Laboratory tests include ceruloplasmin and copper studies (Wilson disease), thyroid function, and drug screening.
\end{itemize}

\section*{Treatment / Management Overview}
\begin{itemize}
  \item Botulinum toxin injections are first-line for focal dystonias (cervical, blepharospasm, laryngeal, limb).
  \item Oral medications include anticholinergics (trihexyphenidyl), baclofen, benzodiazepines, and tetrabenazine.
  \item Deep brain stimulation (pallidal DBS) is highly effective for medication-refractory primary generalized dystonia.
  \item Physical and occupational therapy, splinting, and adaptive equipment optimize function and prevent contractures.
\end{itemize}

\clearpage
